\documentclass[11pt,a4paper]{article}
\usepackage[T1]{fontenc}
\usepackage{lmodern,microtype,amsmath,amssymb,amsthm}
\usepackage[a4paper,margin=25mm,top=22mm,bottom=23mm,headheight=14pt]{geometry}
\usepackage{xcolor,fancyhdr}
\definecolor{deep}{HTML}{174D5F}
\usepackage[colorlinks=true,allcolors=deep,pdftitle={Prime-race oscillation and the Gilbreath triangle}]{hyperref}
\newtheorem{theorem}{Theorem}
\newtheorem{lemma}[theorem]{Lemma}
\newtheorem{corollary}[theorem]{Corollary}
\newcommand{\F}{\mathbb F_2}
\newcommand{\eps}{\varepsilon}
\newcommand{\osc}{\mathcal D}
\newcommand{\dd}{\mathrm D}
\pagestyle{fancy}\fancyhf{}\fancyhead[L]{\small\sffamily Prime-race oscillation and the Gilbreath triangle}
\fancyhead[R]{\small\sffamily Research note}\fancyfoot[C]{\thepage}
\renewcommand{\headrulewidth}{0.3pt}
\setlength{\emergencystretch}{2em}
\raggedbottom
\begin{document}
\thispagestyle{plain}
\noindent{\small\sffamily\bfseries\color{deep} ANALYTIC NUMBER THEORY}\par\medskip
\noindent{\LARGE\bfseries Prime-race oscillation and the Gilbreath triangle}\par\smallskip
\noindent{\large Finite-prefix periods, horizontal blocks, and loss of the largest gaps}\par\smallskip
\noindent{\small 21 November 2026}\par
\begin{abstract}
We ask how closely the running difference between the counts of primes in the two odd residue classes modulo four can follow a line when its slope may be chosen anew for each finite prefix. On the first $N$ odd primes the best maximum error is at least $N^{1/2-\eps}$ for every fixed $\eps>0$ and all sufficiently large $N$. The proof uses opposite prime-power race excursions at two different scales. It gives a lower bound of the same exponent for the least period of each prefix, including a bound allowing symbol changes. In the absolute-difference triangle of normalized prime gaps, a constant-parity block starting in column $j$ at depth $r$ has length $O_\eps((j+r+1)^{2+\eps})$. On column two the residues zero and two modulo four both occur between heights $N^{1/2-\eps}$ and $N$ for every sufficiently large $N$. For the apex of the first $N$ gaps, we exclude the largest two possible values for all sufficiently large $N$, and the largest three when the maximum is divisible by three. The universal bound of one required for the normalized Gilbreath problem remains open.
\end{abstract}
\section{The original target}
Let $p_1=2,p_2=3,\ldots$ be the consecutive primes. Write $(\dd x)_j=|x_{j+1}-x_j|$ and $A_{r,j}=(\dd^r p)_j$, starting at $r=0$. Proth's 1878 question asks whether $A_{r,1}=1$ for every $r\ge1$ \cite{proth,cht}. Let $q_i=p_{i+1}$ be the $i$th odd prime and put
\[
 c_i=\frac{q_{i+1}-q_i}{2}-1,\qquad H_N=\max_{1\le i\le N}c_i,
 \qquad u_N=(\dd^{N-1}c)_1 .
\]
Here $c=(0,0,1,0,1,\ldots)$. Since $A_{1,2}=2$ and $A_{N,2}=2u_N$ for $N\ge2$, the recurrence $A_{r+1,1}=|A_{r,1}-A_{r,2}|$ gives the exact equivalence
\begin{equation}\label{target}
 A_{r,1}=1\ \ (r\ge1)\quad\Longleftrightarrow\quad u_N\le1\ \ (N\ge1).
\end{equation}
We exploit a constraint on the full ordered prime sequence before returning to $u_N$. The distinction between the column index and the depth will be essential.
\newpage
\section{A uniform bound for the prime race with an adjustable slope}
Let $P(t)$ count the odd primes up to $t$. Write $D(t)=\pi(t;4,1)-\pi(t;4,3)$ and let $\chi$ be the nonprincipal character modulo four. Thus $D(t)=\sum_{2<p\le t}\chi(p)$. For $N\ge1$ set $S_m=\sum_{i=1}^m\chi(q_i)$ and $S_0=0$.
\begin{theorem}\label{race}
For every $\eps>0$, for all sufficiently large $X$,
\[
 \inf_{\lambda\in\mathbb R}\max_{2\le t\le X}
         |D(t)-\lambda P(t)|\ \ge\ X^{1/2-\eps}.
\]
In particular $\inf_\lambda\max_{0\le m\le N}|S_m-\lambda m|\ge N^{1/2-\eps}$ for all sufficiently large $N$.
\end{theorem}
\begin{proof}
Define the prime-power race by $Q(t)=\sum_{p^k\le t,\ p>2}\chi(p)^k/k$. Knapowski and Tur\'an's Theorem 2.1 applies unconditionally at modulus four, which appears in their list (1.2) \cite{kt1,kt2}. For every sufficiently large $T$ there are points of both signs in $[T^{1/3},T]$ at which $|Q|$ exceeds $B(T)=T^{1/2-o(1)}$. Their explicit
$B(T)$ is $\sqrt T\exp(-41\log T\,\log_3T/\log_2T)$, where subscripts denote iterated logarithms. Every odd prime square contributes with positive sign. Counting higher powers trivially yields
\begin{equation}\label{powers}
 Q(t)=\lambda P(t)+E(t)+\tfrac12P(\sqrt t)+O(t^{1/3}\log(2t)),
 \qquad E(t)=D(t)-\lambda P(t).
\end{equation}
The error follows by summing $t^{1/k}$ for $3\le k\le \log(t)/\log(2)$.

It suffices to take $0<\eps<1/4$. Suppose $|E(t)|<X^{1/2-\eps}$ throughout $[2,X]$. Fix $1-2\eps<a<1$ and apply the negative excursion at $T=X^a$. At the resulting $t\le T$, both the error in \eqref{powers} and $|E(t)|$ are $o(B(T))$. Dropping the nonnegative square term gives $-\lambda P(t)\ge (1-o(1))B(T)$. Since $P(t)\le T$,
\[
                      -\lambda\ \ge\ T^{-1/2-o(1)}.
\]
At the positive excursion $y\le X$ we now have
$B(X)\le Q(y)\le X^{1/2-\eps}+O(\sqrt y)$, using $\lambda<0$.
Hence $y\ge X^{1-o(1)}$. Substituting again into \eqref{powers} and using $Q(y)>0$ gives $-\lambda P(y)=O(\sqrt y)$. By the prime number theorem $-\lambda\le X^{-1/2+o(1)}$, contradicting the first bound since $a<1$ is fixed. Taking $X=q_N$ identifies the two systems of partial sums and proves the second assertion.
\end{proof}
\noindent
The use of a smaller scale fixes the sign and a lower bound for the magnitude of the slope. The excursion at the larger scale gives an incompatible upper bound. A single large value of the uncentered discrepancy would allow the slope to absorb that value.
\newpage
\section{Finite periods and constant-parity blocks}
A finite word $s_1,\ldots,s_N$ has period $P_0\ge1$ if $s_i=s_{i+P_0}$ whenever both indices exist. The definition permits a partial last period.
\begin{corollary}\label{period}
If the signs of the first $N$ odd primes differ from a $P_0$-periodic sign word at $e$ positions, then, for every $\eps>0$ and all sufficiently large $N$,
\[
                        2(P_0+e)\ \ge\ N^{1/2-\eps}.
\]
In particular their least period is at least $N^{1/2-\eps}$ for all sufficiently large $N$, by first applying the preceding bound with $\eps/2$.
\end{corollary}
\begin{proof}
Extend the periodic word by repeating a period and take its mean $\lambda$. Its partial sums differ from $m\lambda$ by at most $2P_0$ from the last incomplete period. Each substituted sign changes any partial sum by at most two. Apply Theorem~\ref{race}, using a smaller positive $\eps$ to absorb the fixed factor in the last assertion.
\end{proof}
Let $b_i=(q_i-1)/2\pmod 2$, so that $\chi(q_i)=(-1)^{b_i}$. With $S$ denoting the forward shift, let $\delta=1+S$ over $\F$. Absolute and ordinary differencing agree there. In particular
\begin{equation}\label{bin}
 c=1+\delta b\pmod2,\qquad (\dd^r c)\bmod2=\delta^{r+1}b\quad(r\ge1).
\end{equation}
\begin{theorem}\label{strip}
Fix $\eps>0$. Suppose $L$ consecutive entries of $\dd^r c$ starting in column $j\ge1$ have the same parity, with $r\ge0$. Uniformly in $j,r$,
\[
                     L\ \ll_\eps\ (j+r+1)^{2+\eps}.
\]
More generally, if the $L$ entries all lie in a single residue class modulo $d\ge2$, then at least one of the following holds:
\[
 L\ll_\eps(j+r+1)^{2+\eps},\qquad
 d\ll_\eps(j+r+1)^{21/20+\eps}.
\]
For even $d$ the length bound always holds.
\end{theorem}
\begin{proof}
Constancy in \eqref{bin} gives $\delta^{r+2}b=0$ at positions $j,\ldots,j+L-2$, including the case $r=0$. There are $L+r+1$ letters of $b$ in the dependence interval, ending at $K=j+L+r$. For the least power of two $P_0\ge r+2$, $\delta^{P_0}=1+S^{P_0}$. Applying the additional differences on all available positions proves period $P_0$ on the whole interval. If $P_0$ exceeds its length the period condition is vacuous and the length bound already holds. Extend that periodic word to the left. It agrees with the first $K$ prime signs except at most $j-1$ positions. Thus $K^{1/2-\eta}\le 2(P_0+j-1)$ by Corollary~\ref{period}, for each fixed small $\eta>0$ and $K$ large. Choose $\eta$ so $1/(1/2-\eta)<2+\eps$.
\smallskip
For the modulus assertion, write $B=j+r+1$. With any small $\eta>0$, the first $M=\lceil C_\eta B^{2+\eta}\rceil$ entries starting in column $j$ have both parities by the proved length bound, upon increasing $C_\eta$. If $L<M$ we are done. Otherwise two distinct entries there differ by at least $d$. Every entry is bounded by the maximum of its nonnegative top data, drawn from the first $j+r+M$ normalized gaps. Baker, Harman and Pintz give $q_{i+1}-q_i\ll q_i^{21/40}$ \cite{bhp}. Together with $q_k\ll k\log(2k)$ this gives
\[
                   d\ll \big((B+M)\log(2B+2M)\big)^{21/40}.
\]
Take $\eta$ small enough to absorb the logarithm and obtain the stated $21/20+\eps$. A class modulo an even integer fixes parity directly.
\end{proof}
\newpage

\section{The critical second column}
Let $y_0=b_1$ and $y_r=A_{r,2}/2\pmod 2$ for $r\ge1$. For $x_n=b_{n+1}$, $n\ge0$, the identity for iterated differences modulo two gives
\[
 y_r=\sum_{i=0}^r\binom ri x_i\pmod2.
\]
The binomial transform is an involution over $\F$, since
$\sum_{i=k}^n\binom ni\binom ik=\binom nk2^{n-k}$.
\begin{theorem}\label{ray}
Suppose $1\le R\le N$ and the word $y_R,\ldots,y_N$ agrees with a binary sequence of period $P_0$, where $P_0$ is a power of two. For every $\eps>0$ and all sufficiently large $N$,
\[
                      R+P_0\ \gg_\eps N^{1/2-\eps}.
\]
In particular, for $0<\eps<1/2$ and every sufficiently large $N$, the second column contains both residues $0$ and $2$ modulo four at heights between $\lceil N^{1/2-\eps}\rceil$ and $N$.
\end{theorem}
\begin{proof}
Extend the comparison word $v$ periodically back to index zero. Its binomial transform at $n$ equals $(\delta^n v)_0$, which vanishes for $n\ge P_0$ because $\delta^{P_0}v=(1+S^{P_0})v=0$. Thus for $P_0\le n\le N$ inversion gives
\[
 x_n=\sum_{k<R}\binom nk(y_k+v_k)\pmod2 .
\]
Here $\binom nk=0$ for $k>n$. As a function of $n$ the right side is periodic with period the least power of two $J\ge R$. Indeed $(1+z)^{n+J}=(1+z)^n(1+z^J)$ in $\F[z]$, so all coefficients below $J$ agree. Hence among the first $N+1$ prime signs at most $P_0$ must be changed to obtain a $J$-periodic word. Corollary~\ref{period} and $J<2R$ prove the result. For the last assertion use period one and apply the bound with $\eps/2$.
\end{proof}
For example, a value congruent to $2$ modulo four can equal either $2$ or $6$. The former preserves a western value of one. The latter takes it to five. The next estimate distinguishes actual magnitudes.
\newpage
\section{Extinction of the two highest values on the anchored cone}
The cone of a cell is the triangle generated by the consecutive input entries on which the cell depends. An elementary fact behind Chase--Hunter--Tao's parentage lemma \cite[Lemma 3.8]{cht} will be used twice.
\begin{lemma}\label{max}
If all top entries are in $[0,H]$, $H>0$, and the apex is $H$, the entire top word is contained in $\{0,H\}$.
\end{lemma}
\begin{proof}
A difference of $H$ has parents $0,H$. In a row with entries $0,H$ containing $H$, start at this pair of parents. A prescribed difference of $0$ or $H$ and one known extreme parent force the other parent to be extreme. Propagating in both directions recovers an entire parent row in $\{0,H\}$ which contains $H$. Repeat up to the top.
\end{proof}
\begin{lemma}\label{mod3}
A block of consecutive odd primes with gaps in $\{u,v\}$, $u,v\le K$ and $3\nmid\gcd(u,v)$, has $O(\log(2K))$ gaps, uniformly in its location.
\end{lemma}
\begin{proof}
Omit an initial $3$ if present. Both allowed prime residues modulo three are nonzero. If the two gap residues coincide and are nonzero, two successive steps would encounter zero. If exactly one is nonzero, that gap can occur at most once, leaving at most two constant-gap runs. If the residues are opposite and nonzero, the gap types must alternate. In the last case each parity of the prime positions gives an arithmetic progression of step $u+v$.

For a progression of primes of positive even step $h$, let $\ell$ be the least prime not dividing $h$. Any $2\ell$ terms include two distinct multiples of $\ell$, which cannot both be prime. The product of the primes below $\ell$ divides $h$, so the Chebyshev bound $\sum_{p\le x}\log p\gg x$ implies $\ell\ll\log(2h)$. This bounds each progression in the preceding cases.
\end{proof}
\noindent
It follows that in any finite triangle of normalized odd-prime gaps bounded by $H\ge2$, every cell of value $H$ has depth $O(\log(2H))$. Indeed its top cone has gap sizes $2$ and $2(H+1)$ by Lemma~\ref{max}, and its length is one more than its depth.

\begin{corollary}\label{depth}
Consider the triangle on $c_1,\ldots,c_N$ with maximum $H_N$. If a full row at depth $t$ is constant modulo $d$, where $2\le d\le H_N$, then
\[
                        d\ll_\eps (t+2)^{21/20+\eps}.
\]
Thus a full row contained in $\{0,d\}$ and attaining $d$ can occur only at depth $t+2\gg_\eps d^{20/21-\eps}$.
\end{corollary}
\begin{proof}
Apply Theorem~\ref{strip} with $j=1$ and $L=N-t$. Its modulus bound suffices directly. In the remaining case $N\ll_\eta(t+2)^{2+\eta}$. Then $d\le H_N\ll(N\log(2N))^{21/40}$ by the same prime-gap estimate. Choose $\eta>0$ sufficiently small.
\end{proof}
\newpage
\begin{theorem}\label{contraction}
For every sufficiently large $N$, $u_N\le H_N-2$. For those $N$ with $3\mid H_N$, one also has $u_N\le H_N-3$.
\end{theorem}
\begin{proof}
Write $H=H_N$. Then $H\to\infty$ since arbitrarily long strings of consecutive composites occur. Take $T=\lceil C\log(2H)\rceil$ with a sufficiently large absolute $C$. Every entry at depth $T$ is at most $H-1$. If $u_N=H-1$, Lemma~\ref{max} applied from that row down requires the entire row to be contained in $\{0,H-1\}$.

The following argument works with either $h=H-1$ or $h=H-2$. Put $M=\lfloor\sqrt H\rfloor$. Because $H\le q_{N+1}\ll N\log(2N)$, we have $T<M\le N$ for all sufficiently large $N$. The first $M$ input values are at most $q_{M+1}\ll M\log(2M)<h$. All entries of their subtriangle are therefore less than $h$. If the whole row $T$ belonged to $\{0,h\}$, its first $M-T$ entries would be zero. Theorem~\ref{strip} bounds this length by $O((T+2)^3)$, contradicting $M-T\asymp\sqrt H$. This excludes $h=H-1$.

Now let $3\mid H$. No entry of value $H$ is at positive depth. Otherwise the parent rows would give a top block in $\{0,H\}$ of length at least two containing $H$. Both allowed prime gaps $2,2(H+1)$ are $2$ modulo three. No two consecutive such gaps can start at a prime greater than three. The exceptional pair starting at three consists of gaps $2,2$ and does not contain $H$ in its normalized word.

The first difference row is consequently bounded by $H-1$. If a cell of value $H-1$ occurs at depth $t\ge1$, apply Lemma~\ref{max} from that first row to the cell. Its first-row ancestors belong to $\{0,H-1\}$ and include $H-1$. The original normalized inputs above them are all congruent modulo $H-1$. They lie in $[0,H]$ and attain two values differing by $H-1$, so their alphabet is either $\{0,H-1\}$ or $\{1,H\}$. The respective ordinary gap pairs are $\{2,2H\}$ and $\{4,2(H+1)\}$. Lemma~\ref{mod3} applies in both cases, giving $t=O(\log(2H))$. Increasing $C$ ensures every entry of row $T$ is at most $H-2$. An apex $H-2$ would now force a row in $\{0,H-2\}$ and is excluded by the preceding small-prefix argument.
\end{proof}
\begin{corollary}\label{wbound}
Let $G_N=\max_{i\le N}(q_{i+1}-q_i)=2(H_N+1)$. In the original prime triangle, for all sufficiently large $N$,
\[
 A_{N+1,1}\le G_N-7,\qquad
  A_{N+1,1}\le G_N-9\quad\text{if }3\mid H_N.
\]
\end{corollary}
\begin{proof}
Put $w_r=A_{r,1}$ and $B_r=A_{r,2}$. Since $w_r$ is positive and odd and $B_r$ is even, $|w_r-B_r|\le\max(w_r,B_r-1)$. Iterate from $w_1=1$. For large $i$, $B_i=2u_i\le2H_i-4\le2H_N-4$ for $i\le N$. The finitely many earlier values are absorbed as $H_N\to\infty$. If $3\mid H_N$, then either $H_i\le H_N-1$, or $H_i=H_N$ and the stronger bound of Theorem~\ref{contraction} applies. In both cases $B_i\le2H_N-6$ for all sufficiently large $i\le N$.
\end{proof}
\section{Where the argument stops}
A bound $u_N\le H_N-2$ still permits $u_N=2$ for arbitrarily large $N$. An argument that simply repeats the logarithmic extinction step in every translated window faces the following concrete obstruction. For $H\ge2$ with $3\nmid H$, the normalized word $(H,0,1,\ldots,H)$ produces $(H,1,\ldots,1)$ and then $(H-1,0,\ldots,0)$. Its left entry remains $H-1$ at every depth from $2$ through $H+1$.
\newpage
The corresponding integer offsets of the prospective primes are
\[
 \{0\}\ \cup\ \{2(H+1)+j(j+1):0\le j\le H+1\}.
\]
They omit a residue class modulo every prime. All offsets are even. Modulo three the quadratic takes the values $0,2$ and the translated pair contains zero. Modulo an odd prime $\ell\ge5$ the quadratic takes at most $(\ell+1)/2$ values because $j$ and $-1-j$ agree, so adding zero still leaves a missing class. Under Dickson's prime-tuples conjecture the indicated pattern would occur among consecutive primes for each such $H$. Indeed assign a distinct prime larger than the span to each omitted integer offset and fix the translate by the Chinese remainder theorem to make all those positions composite. The linear forms at designated offsets remain admissible, since none meets a forced divisor and the other local tests are unchanged. This conditional example explains a limitation of the repeated extinction strategy. All bounds in Theorems~\ref{race}, \ref{strip} and \ref{contraction} were proved unconditionally.

Chase, Hunter and Tao's deterministic sufficient criterion \cite[Thm.~1.6]{cht} controls initial heights, zero blocks throughout the array, and shallow two-valued blocks which may start far to the right. Theorem~\ref{strip} is informative for blocks near the left boundary. Its dependence on the starting column leaves the right-half obstruction untreated at the required scales. Theorem~\ref{contraction} concerns the fixed apex in \eqref{target}. It gives neither the universal bound $u_N\le1$ nor infinitely many ones on the original western edge \cite{bcz}.

\paragraph{Relation to earlier work.}
The corpus manuscript \cite{vm} already uses modular restrictions and primorial divisibility for bounded-depth obstruction words. Qualitative nonperiodicity of prime residues and the role of prime squares in it were discussed in 2014 \cite{mo}. Our estimate concerns every sufficiently long finite prefix and allows the fitted slope to vary with the prefix. The new use within this argument is to force short periods from a hypothetical surviving maximum, then contradict them by opposite race excursions at different scales. The originality of these statements has not been independently established. The historical problem and its quantifiers are retained. Its complete resolution and the priority criterion remain unmet.

\begin{thebibliography}{10}\small
\bibitem{proth} F. Proth, Sur la s\'erie des nombres premiers, \emph{Nouvelle Correspondance Math\'ematique} 4 (1878), 236--240. \href{https://gdz.sub.uni-goettingen.de/download/pdf/PPN598948236_0004/LOG_0088.pdf}{Original paper}.
\bibitem{kt1} S. Knapowski and P. Tur\'an, Comparative prime-number theory I, \emph{Acta Math. Acad. Sci. Hungar.} 13 (1962), 299--314, eq.~(3.6). \href{https://personal.math.ubc.ca/~gerg/teaching/592-Fall2018/papers/1962.Knapowski_1.pdf}{Paper}.
\bibitem{kt2} S. Knapowski and P. Tur\'an, Comparative prime-number theory II, \emph{ibid.} 13 (1962), 315--342, list (1.2) and Theorem 2.1. \href{https://personal.math.ubc.ca/~gerg/teaching/592-Fall2018/papers/1962.Knapowski_2.pdf}{Paper}.
\bibitem{bhp} R. C. Baker, G. Harman and J. Pintz, The difference between consecutive primes, II, \emph{Proc. London Math. Soc.} 83 (2001), 532--562. \href{https://doi.org/10.1112/plms/83.3.532}{DOI}.
\bibitem{cht} Z. Chase, Z. Hunter and T. Tao, Gilbreath's conjecture: a Cram\'er random model and a deterministic analysis (2026). \href{https://arxiv.org/abs/2607.08712}{arXiv:2607.08712}.
\bibitem{bcz} R. N. Bhat, C. Cobeli and A. Zaharescu, Filtered rays over iterated absolute differences on layers of integers (2023). \href{https://arxiv.org/abs/2309.03922}{arXiv:2309.03922}.
\bibitem{mo} solaris and Lucia, The prime numbers modulo $k$, are not periodic, MathOverflow (2014). \href{https://mathoverflow.net/questions/168378}{Question and answers}.
\bibitem{vm} Keio, Integrated Structures and Boundary Spectra in Sparse Prime Codings, VMRC 003, section on fixed-depth Gilbreath obstruction languages. \href{https://github.com/QEEDDD/Verifiable-Mathematical-Reasoning-Corpus/tree/c2ac34ea4622778187173f41bb1996deee8081b7}{Repository revision c2ac34ea4622}.
\end{thebibliography}
\end{document}
